A pair table appears to be one of the simplest objects in sequence modeling.  Choose two positions,
list the possible residue identities along the rows and columns, and attach a number to every
combination.  The result is finite, visible, and easy to store.  Yet before it can represent an
interaction, a basic ambiguity must be resolved: which part of the table belongs specifically to
the pair?

Suppose a table $T(a,b)$ is changed by
\begin{equation}
T(a,b)\longmapsto T(a,b)+u(a)+v(b)+c.
\label{eq:book-table-gauge-opening}
\end{equation}
In a conditional model, the added row term may be absorbed into a local field at the target; in a
joint energy, row and column terms can be redistributed among one-body potentials.  Different raw
tables can therefore encode the same observable distribution or the same categorical response.
Taking the matrix at face value confuses a parameterization with an identified pair object.

The remedy is often called gauge fixing, but the phrase can make the operation sound like a
convention chosen for numerical convenience.  The construction in this chapter is stronger.  A
reference distribution at each endpoint defines which directions count as constants and one-body
effects.  Double projection removes exactly those directions.  What remains is the canonical
bivariate component relative to that reference.

The geometric interpretation is the key.  Position $i$ carries a $(q-1)$-dimensional space of
centered categorical contrasts behind its $q$ labels.  A pair interaction is a map from the source
contrast space to the target contrast space.  The resulting protein graph is a bundle of local
categorical spaces connected by typed linear operators; scalar edge weights are compressed shadows
of these maps.

This shift resolves several problems at once.  It explains why raw Potts tables are not directly
comparable across gauges, why reference changes require a transport law, why scalar norms lose
substitution direction, and why directed model responses can be compared with static couplings
without pretending that their original parameters are the same.  The formal definitions below
make these claims exact.
